% ============================================================
% Round 01: What Does Image Editing Post-Training Optimize?
% ============================================================

\subsection{Round 01：图像编辑后训练到底在优化什么}

\subsubsection{本轮问题}

\begin{conceptbox}
\textbf{学习目标：} 先不急着上 DDPO、Flow-GRPO 或 HP-Edit。先弄清楚图像编辑后训练的基本对象：输入是什么、输出是什么、什么叫“编辑成功”、reward 为什么难设计、为什么 SFT 不够。\\
\textbf{主线位置：}
\[
\text{图像编辑任务本体}
\to
\text{Diffusion / Flow 后训练基础}
\to
\text{Reward / Preference}
\to
\text{Diffusion-DPO / DDPO / Flow-GRPO / HP-Edit}
\]
\end{conceptbox}

\subsubsection{图像编辑和 Text-to-Image 的区别}

普通 text-to-image 生成的输入通常是：

\[
c = \text{文本 prompt}
\]

输出是：

\[
I = \text{生成图像}
\]

目标大概是：

\[
\text{图像符合 prompt，且质量高}
\]

但图像编辑的输入是：

\[
(I_\text{src}, c)
\]

其中：

\begin{itemize}[leftmargin=1.5em]
  \item $I_\text{src}$：原始图像；
  \item $c$：编辑指令；
  \item $I_\text{edit}$：编辑后的输出图像。
\end{itemize}

目标不只是“生成好图”，而是：

\[
\text{在保留原图大部分内容的同时，精确完成指定编辑}
\]

\begin{center}
{\small
\begin{tabular}{p{3cm}p{4.8cm}p{4.8cm}}
\hline
 & \textbf{Text-to-Image} & \textbf{Image Editing} \\
\hline
输入
&
文本 prompt
&
原图 + 编辑指令
\\
输出
&
新图像
&
编辑后的图像
\\
核心约束
&
符合文本、质量高
&
该改的要改，不该改的要保持
\\
失败类型
&
不符合 prompt、质量差
&
没改对、改多了、身份漂移、背景乱变、质量下降
\\
\hline
\end{tabular}
}
\end{center}

\begin{warnbox}
\textbf{关键区别：} 图像编辑比 text-to-image 多了一个强约束：输出必须和输入图像保持关系。不能只看输出图像本身好不好。
\end{warnbox}

\subsubsection{图像编辑的四个目标}

一个好的图像编辑结果通常要同时满足四件事：

\[
R
=
R_\text{instruction}
+
R_\text{preservation}
+
R_\text{quality}
+
R_\text{preference}
\]

分别是：

\begin{enumerate}[leftmargin=1.5em]
  \item \textbf{Instruction following：} 是否完成了编辑指令。
  \item \textbf{Preservation：} 不该改的区域是否保持。
  \item \textbf{Visual quality：} 画面是否自然、清晰、无伪影。
  \item \textbf{Human preference：} 人类是否更喜欢这个结果。
\end{enumerate}

举例：

\begin{center}
{\small
\begin{tabular}{p{3cm}p{9.8cm}}
\hline
\textbf{编辑指令} & 把照片里的红色车改成蓝色车 \\
\hline
Instruction & 车是否真的变成蓝色 \\
Preservation & 背景、车形状、光照、视角是否保持 \\
Quality & 蓝色是否自然，有没有涂抹、边缘破损 \\
Preference & 人类是否觉得这个结果更真实、更符合意图 \\
\hline
\end{tabular}
}
\end{center}

\begin{intubox}
\textbf{直觉：} 图像编辑不是“越改越明显越好”，而是“只改用户要求的部分，而且改得自然”。
\end{intubox}

\subsubsection{为什么 Reward 更难设计}

文本 RLVR 里，数学题可以用：

\[
R = \mathbb{1}[\text{final answer correct}]
\]

但图像编辑很少有这么简单的 verifier。

比如指令是：

\[
\text{把天空变成黄昏}
\]

我们要判断：

\begin{itemize}[leftmargin=1.5em]
  \item 天空有没有变成黄昏？
  \item 房子、人物、道路有没有不该变？
  \item 光照变化是否自然？
  \item 颜色是否符合人类偏好？
  \item 如果原图没有天空怎么办？
\end{itemize}

这些都不是简单的 0/1 检查。

所以图像编辑 reward 通常是复合的：

\[
R(I_\text{src},c,I_\text{edit})
=
\lambda_1 R_\text{edit}
+
\lambda_2 R_\text{preserve}
+
\lambda_3 R_\text{quality}
+
\lambda_4 R_\text{pref}
\]

其中：

\begin{itemize}[leftmargin=1.5em]
  \item $R_\text{edit}$：编辑是否符合指令；
  \item $R_\text{preserve}$：输入和输出中不该变的内容是否保持；
  \item $R_\text{quality}$：图像质量 / 美学；
  \item $R_\text{pref}$：人类偏好模型或 VLM scorer 的评分。
\end{itemize}

\begin{warnbox}
\textbf{难点：} 这些 reward 经常互相冲突。提高编辑强度可能破坏保真度；提高保真度可能导致编辑不明显；提高美学可能让模型擅自改背景。
\end{warnbox}

\subsubsection{图像编辑里的 Policy、Action、Trajectory}

在文本侧，我们有：

\[
\pi_\theta(y_t\mid x,y_{<t})
\]

它每一步 action 是下一个 token。

在图像生成 / 编辑侧，模型不是一个 token 一个 token 生成，而是在 latent 或 pixel 空间逐步去噪 / 积分。

可以粗略写成：

\[
z_T
\to
z_{T-1}
\to
\cdots
\to
z_0
\to
I_\text{edit}
\]

其中：

\begin{itemize}[leftmargin=1.5em]
  \item $z_T$：噪声 latent；
  \item $z_t$：中间 latent 状态；
  \item $z_0$：最终 clean latent；
  \item $I_\text{edit}$：VAE decoder 解码后的图像。
\end{itemize}

如果是图像编辑，条件里还包括：

\[
I_\text{src}, c
\]

所以生成过程可以看成：

\[
p_\theta(I_\text{edit}\mid I_\text{src}, c)
\]

或者更细：

\[
p_\theta(z_{t-1}\mid z_t, I_\text{src}, c, t)
\]

\begin{center}
{\small
\begin{tabular}{p{3cm}p{4.8cm}p{4.8cm}}
\hline
\textbf{概念} & \textbf{文本侧} & \textbf{图像编辑侧} \\
\hline
State
&
prompt + 已生成 tokens
&
当前 latent $z_t$ + 原图 + 编辑指令
\\
Action
&
下一个 token
&
去噪一步 / flow step / latent 更新
\\
Trajectory
&
完整回答 token 序列
&
完整 denoising / flow 轨迹
\\
Reward
&
回答质量 / 正确性
&
最终编辑图像质量、指令遵循、保真度
\\
\hline
\end{tabular}
}
\end{center}

\begin{summarybox}
\textbf{最小映射：} 文本 RL 的 trajectory 是 token 序列；图像编辑 RL 的 trajectory 是从噪声 latent 到最终编辑图像的生成过程。
\end{summarybox}

\subsubsection{为什么 SFT 不够}

图像编辑 SFT 通常使用 paired editing data：

\[
(I_\text{src}, c, I_\text{target})
\]

训练模型模仿目标图像：

\[
\text{给定原图和编辑指令，生成目标编辑图}
\]

这对应文本侧的 SFT：

\[
(x,y^\ast)
\]

但 SFT 有几个限制：

\begin{enumerate}[leftmargin=1.5em]
  \item \textbf{目标图像不一定唯一。} 同一个指令可能有很多合理编辑结果。
  \item \textbf{配对数据难且贵。} 高质量真实编辑前后图很难大规模获取。
  \item \textbf{像素级模仿不等于人类偏好。} 两张都符合指令的图，人类可能明显更喜欢其中一张。
  \item \textbf{难处理复合目标。} 编辑正确、区域保持、审美质量经常互相冲突。
\end{enumerate}

所以后训练要补上的东西是：

\[
\text{不只是模仿目标图，而是优化人类更喜欢的编辑结果}
\]

\begin{summarybox}
\textbf{一句话：} 图像编辑 SFT 教模型“怎么做编辑”；偏好 / RL 后训练教模型“哪种编辑结果更好”。
\end{summarybox}

\subsubsection{偏好数据在图像编辑里长什么样}

文本 DPO 的数据是：

\[
(x,y_w,y_l)
\]

图像编辑偏好数据变成：

\[
(I_\text{src}, c, I_w, I_l)
\]

其中：

\begin{itemize}[leftmargin=1.5em]
  \item $I_w$：人类更喜欢的编辑结果；
  \item $I_l$：人类不喜欢的编辑结果；
  \item 两者都对应同一个原图和同一个编辑指令。
\end{itemize}

这就可以训练：

\begin{itemize}[leftmargin=1.5em]
  \item 图像 reward model / scorer；
  \item Diffusion-DPO 类 preference objective；
  \item RLHF / RLVR 风格的在线优化。
\end{itemize}

\begin{intubox}
\textbf{直觉：} 图像编辑偏好比较必须是条件化的。不能只问“哪张图更好看”，而要问“在这张原图和这条编辑指令下，哪张编辑结果更好”。
\end{intubox}

\subsubsection{图像编辑后训练的三类路线}

\begin{center}
{\small
\begin{tabular}{p{3cm}p{4.8cm}p{4.8cm}}
\hline
\textbf{路线} & \textbf{数据 / 信号} & \textbf{代表方法} \\
\hline
SFT
&
$(I_\text{src}, c, I_\text{target})$
&
instruction editing fine-tuning
\\
Preference Optimization
&
$(I_\text{src}, c, I_w, I_l)$
&
Diffusion-DPO, preference alignment
\\
RL / RLHF
&
模型在线生成 $I_\text{edit}$，scorer 打 reward
&
DDPO, Flow-GRPO, HP-Edit
\\
\hline
\end{tabular}
}
\end{center}

三者关系可以写成：

\[
\text{SFT}
\to
\text{Preference}
\to
\text{Online RL}
\]

和文本侧非常像：

\[
\text{SFT}
\to
\text{DPO}
\to
\text{GRPO / RLVR}
\]

但区别是图像侧的 policy 和 trajectory 不再是 token，而是 diffusion / flow 生成过程。

\subsubsection{本轮最小结论}

\begin{summarybox}
\textbf{图像编辑后训练到底在优化什么？}\\
给定原图 $I_\text{src}$ 和编辑指令 $c$，优化模型生成的编辑结果 $I_\text{edit}$，使它同时满足：
\[
\text{指令遵循}
+
\text{原图保持}
+
\text{视觉质量}
+
\text{人类偏好}
\]
核心难点是：reward 是复合且冲突的，trajectory 是 denoising / flow 过程，最终图像 reward 需要分配给中间生成步骤。
\end{summarybox}

\subsubsection{本轮检查题}

\begin{enumerate}[leftmargin=1.5em]
  \item 图像编辑和 text-to-image 生成最大的区别是什么？
  \item 为什么图像编辑不能只优化美学分数？
  \item 图像编辑里的 policy、action、trajectory 分别是什么？
  \item 为什么图像编辑 SFT 不够？
  \item 图像编辑偏好数据为什么必须包含原图和编辑指令？
\end{enumerate}
